\thispagestyle{empty}\vspace*{\fill}
\begin{minipage}{128mm}\fontsize{9}{13}\selectfont\color{otMuted}
{\otSemi\color{otInk}Mathématiques · Troisième\par Édition 2028-2029 · Nouveau programme · Brevet 2029}\par\medskip
Conception et mise en page : collection OptiTAB.\par
270 exercices de chapitre, 18 défis et deux brevets blancs originaux, corrigés.\par
Charte graphique de la collection OptiTAB.\par\bigskip
Édition du 23 octobre 2026 · France · Année scolaire 2028-2029.\par\bigskip
Cette édition suit le programme de mathématiques du cycle 4 fixé par l’arrêté
du 18 février 2026 (BO n° 10 du 5 mars 2026), applicable en Troisième à la
rentrée 2028. Elle prépare la session 2029 du diplôme national du brevet.\par\bigskip
Au 23 octobre 2026, aucun texte propre à la session 2029 n’était publié. Les
brevets blancs suivent la définition de l’épreuve de mathématiques en vigueur
à compter de la session 2027 (note de service du 11 septembre 2026, BO spécial
n° 4 du 17 septembre 2026). Vérifier avant l’examen qu’aucun texte plus récent
ne l’a modifiée.\par\bigskip
L’ordre des chapitres, les exercices et les barèmes d’entraînement sont des
choix éditoriaux. Les défis sont facultatifs. Les situations et données
numériques sont construites pour l’apprentissage ; elles ne constituent pas
des données économiques ou scientifiques relevées.\par\bigskip
Les schémas de principe ne permettent pas de mesurer une longueur inconnue.
Les coordonnées et graduations explicites font foi pour les graphiques.
Sauf demande d’arrondi, conserver les résultats exacts.\par\bigskip
© OptiTAB, 2026.\par\otSite
\end{minipage}\clearpage
